\begin{frame}{엔트로피, 교차 엔트로피, KL 발산}
  \begin{itemize}
    \item 정보량의 \textbf{기댓값} = \kb{엔트로피}(entropy):
      \[H(p) = -\sum_k p_k \log_2 p_k\]
      분포 $p$ 를 따르는 사건들을 평균적으로 관찰할 때 얻는
      정보량이자, 이 분포를 \textbf{최적 부호화(encoding)} 하는 데
      필요한 평균 비트 수 (Week 7.2 결정트리에서 ``가장 잘 나누는
      질문'' 기준으로 다시 만남).
    \item 실제 분포는 $p$ 인데, (틀린) 다른 분포 $q$ 를 믿고
      부호화하면 필요한 \textbf{평균 비트 수} = \kb{교차 엔트로피}:
      \[H(p, q) = -\sum_k p_k \log_2 q_k\]
    \item $q = p$ 일 때 $H(p, q) = H(p)$ 로 \textbf{최솟값}.
      $q$ 가 $p$ 에서 멀어질수록 커지고, 그 초과분
      $H(p, q) - H(p)$ = \kb{KL 발산}($D_{KL}(p\|q)$)
      (Week 15 VAE의 ELBO에서 다시 만남).
  \end{itemize}
\end{frame}
